%% kapa: exam=2 round=34 year=2023 period=3 subject=compensation_law subject_ko=감정평가 및 보상법규 seq=3 number=3 points=20
%% source: Q-Net 감정평가사 기출문제 (cst00309) artlSeq=5212168 file=raw/1008/5212168/3교시 23년도 감정평가사 2차시험_감정평가 및 보상법규.pdf
%% conversion: kordoc 4.24.51 -> tools/gen2.py (auto); review=proofread-gpt-sol; flags=none
\begin{longquestion}{【문제 3】}
A감정평가법인(이하 ‘A법인’이라 함)에 근무하는 B감정평가사(이하 ‘B’라 함)는 2020. 4.경 甲소유의 토지(이하 ‘甲토지’라 함)를 감정평가하면서 甲토지와 이용가치가 비슷하다고 인정되는 ｢부동산 가격공시에 관한 법률｣에 따른 표준지공시지가를 기준으로 감정평가를 하지도 않았고 적정한 실거래가보다 3배 이상 차이가 나는 금액으로 甲토지를 감정평가하였다. 그러나 그 사실은 3년여가 지난 후 발견되었고 이에 따라 국토교통부장관은 감정평가관리ㆍ징계위원회(이하 ‘위원회’라 함)에 징계의결을 요구하였으며 위원회는 3개월의 업무정지를 의결하였고, 국토교통부장관은 위원회의 의결에 따라 2023. 7. 10. B에 대해서 3개월의 업무정지처분(2023. 8. 1.부터)을 결정하였으며 A법인과 B에게 2023. 7. 10. 위 징계사실을 통보하였다. 이에 B는 위 징계가 위법하다는 이유로 2023. 7. 14. 취소소송을 제기하면서 집행정지를 신청하였다. 집행정지의 인용가능성과 본안에서 B의 청구가 기각되는 경우 징계의 효력과 국토교통부장관이 취해야 할 조치에 관하여 설명하시오. (20점) \par
\begin{framedtext}
참조 조문 \par
<감정평가 및 감정평가사에 관한 법률> \par
제39조(징계) ① 국토교통부장관은 감정평가사가 다음 각 호의 어느 하나에 해당하는 경우에는 제40조에 따른 감정평가관리ㆍ징계위원회의 의결에 따라 제2항 각 호의 어느 하나에 해당하는 징계를 할 수 있다. 다만, 제2항제1호에 따른 징계는 제11호, 제12호를 위반한 경우 및 제27조를 위반하여 다른 사람에게 자격증ㆍ등록증 또는 인가증을 양도 또는 대여한 경우에만 할 수 있다. \par
1. 제3조제1항을 위반하여 감정평가를 한 경우 \par
2. 제3조제3항에 따른 원칙과 기준을 위반하여 감정평가를 한 경우 \par
<생략> \par
⑦ 제1항에 따른 징계의결은 국토교통부장관의 요구에 따라 하며, 징계의결의 요구는 위반사유가 발생한 날부터 5년이 지나면 할 수 없다. \par
제39조의2(징계의 공고) ① 국토교통부장관은 제39조제1항 및 제2항에 따라 징계를 한 때에는 지체 없이 그 구체적인 사유를 해당 감정평가사, 감정평가법인등 및 협회에 각각 알리고, 그 내용을 대통령령으로 정하는 바에 따라 관보 또는 인터넷 홈페이지 등에 게시 또는 공고하여야 한다. \par
제40조(감정평가관리ㆍ징계위원회) ① 다음 각 호의 사항을 심의 또는 의결하기 위하여 국토교통부에 감정평가관리ㆍ징계위원회(이하 “위원회”라 한다)를 둔다. \par
<생략> \par
4. 제39조에 따른 징계에 관한 사항 \par
<감정평가 및 감정평가사에 관한 법률 시행령> \par
제29조(인가취소 등의 기준) 법 제32조제1항에 따른 감정평가업자의 설립인가 취소와 업무정지의 기준은 별표 3과 같다. \par
[별표 3] 감정평가업자의 설립인가 취소와 업무정지의 기준(제29조 관련) \par
1. 일반기준 \par
가. 위반행위의 횟수에 따른 행정처분의 기준은 최근 1년간(제2호하목의 경우에는 최근 3년간을 말한다) 같은 위반행위(근거 법조문 내에서 위반행위가 구분되어 있는 경우에는 그 구분된 위반행위를 말한다)로 행정처분을 받은 경우에 적용한다. 이 경우 위반횟수는 같은 위반행위에 대하여 행정처분을 받은 날과 그 처분 후에 다시 같은 위반행위를 하여 적발된 날을 각각 기준으로 하여 계산한다. \par
<생략> \par
다. 국토교통부장관은 위반행위의 동기ㆍ내용 및 위반의 정도 등을 고려하여 처분기준의 2분의 1 범위에서 그 기간을 늘릴 수 있다. 다만, 늘리는 경우에도 총 업무정지기간은 2년을 넘을 수 없다. \par
2. 개별기준 {\small\centering
\resizebox{\linewidth}{!}{%
\begin{tabular}{|p{6.0cm}|c|c|c|c|}\hline
\multirow{2}{*}{위반행위} & \multirow{2}{*}{\makecell{근거\\법조문}} & \multicolumn{3}{c|}{행정처분기준} \\\cline{3-5}
 & & 1차 위반 & 2차 위반 & \makecell{3차 이상\\위반} \\\hline
라. 법 제3조제1항을 위반하여 감정평가를 한 경우 & \makecell{법 제32조\\제1항제4호} & \makecell{업무정지\\1개월} & \makecell{업무정지\\3개월} & \makecell{업무정지\\6개월} \\\hline
마. 법 제3조제3항에 따른 원칙과 기준을 위반하여 감정평가를 한 경우 & \makecell{법 제32조\\제1항제5호} & \makecell{업무정지\\1개월} & \makecell{업무정지\\2개월} & \makecell{업무정지\\4개월} \\\hline
\end{tabular}}\par}
\end{framedtext}
\end{longquestion}
